\documentclass[11pt]{article}
\usepackage[margin=1in]{geometry}
\usepackage{amsmath, amssymb}

\title{Exercise 3.19: $q_\pi(s,a)$ from the Backup Diagram}
\author{Sutton \& Barto, Section 3.5}
\date{}

\begin{document}
\maketitle

\section*{Problem}

The value of an action, $q_\pi(s, a)$, depends on the expected next reward and
the expected sum of the remaining rewards. The backup diagram is rooted at a
state--action pair $(s, a)$ and branches to possible next states.

Give the equation for $q_\pi(s, a)$ in terms of $R_{t+1}$ and $v_\pi(S_{t+1})$,
given $S_t = s$ and $A_t = a$. First with an expectation (not conditioned on
following the policy), then written out explicitly using $p(s', r \mid s, a)$
with no expected-value notation.

\section*{Solution}

\subsection*{With expectation notation}

\[
    q_\pi(s, a) = \mathbb{E}\!\left[\, R_{t+1} + \gamma\, v_\pi(S_{t+1})
    \mid S_t = s,\, A_t = a \,\right]
\]

Note: the expectation is \emph{not} conditioned on following $\pi$. Once we
condition on the action $a$, the next state and reward are determined entirely
by the dynamics $p$, not by the policy.

\subsection*{Written out explicitly}

\[
    \boxed{q_\pi(s, a) = \sum_{s'} \sum_{r} p(s', r \mid s, a)
    \left[\, r + \gamma\, v_\pi(s') \,\right]}
\]

This corresponds to the lower half of the $v_\pi$ backup diagram (Figure~3.4,
left): from the action node, we branch to each possible next state $s'$ and
reward $r$ with probability $p(s', r \mid s, a)$, and read off $v_\pi(s')$
at the leaves.

\subsection*{Relationship to the full Bellman equation}

Substituting the result of Exercise~3.18 into this equation (or vice versa)
recovers the full Bellman equation~(3.14):
\[
    v_\pi(s) = \sum_{a} \pi(a \mid s) \sum_{s'} \sum_{r}
    p(s', r \mid s, a) \left[\, r + \gamma\, v_\pi(s') \,\right]
\]

Similarly, substituting this equation into the result of Exercise~3.17 gives
the Bellman equation for $q_\pi$ purely in terms of itself.

\end{document}
